\uuid{18sl}
\titre{Primitives usuelles}
\auteur{hanani ; mzari}
\organisation{exo7}
\licence{CC BY-NC-SA 4.0}
\datecreate{2020-01-20}
\niveau{L1}
\module{Analyse}
\chapitre{Calcul d'intégrales}
\sousChapitre{Autre}
\difficulte{1}
\theme{Calculs d'intégrales, primitives usuelles}
\exo7qcmid{lille-2.integrales.002}

\contenu{
\begin{question}
\qclassification{127.99}
\qtags{motcle=primitives usuelles}
Parmi les égalités suivantes, cocher celles qui sont vraies :
\begin{answers}
\bad{$\displaystyle \int _0^1\mathrm{e}^{x}\,\mathrm{d}x=\mathrm{e}-1$ et $\displaystyle \int _0^1\mathrm{e}^{2x}\,\mathrm{d}x=\mathrm{e}^2-1$.}
\good{$\displaystyle \int _0^{\pi/4}\sin (2x)\,\mathrm{d}x=\frac{1}{2}$ et $\displaystyle \int _0^{\pi/4}\cos (2x)\,\mathrm{d}x=\frac{1}{2}$.}
\bad{$\displaystyle \int _0^1\frac{\mathrm{d}x}{x+1}=\ln 2$ et $\displaystyle \int _0^1\frac{\mathrm{d}x}{x+2}=\ln 3$.}
\good{$\displaystyle \int _0^1\frac{\mathrm{d}x}{(1+x)^2}=\frac{1}{2}$ et $\displaystyle \int _0^1\frac{\mathrm{d}x}{1+x^2}=\frac{\pi}{4}$.}
\end{answers}
\vskip2mm
\begin{explanations}
On a : $\displaystyle \int _0^1\mathrm{e}^{x}\,\mathrm{d}x=\Big[\mathrm{e}^x\Big]_0^1=\mathrm{e}-1$ mais $\displaystyle \int _0^1\mathrm{e}^{2x}\,\mathrm{d}x=\left[\frac{\mathrm{e}^{2x}}{2}\right]_0^1=\frac{\mathrm{e}^2-1}{2}$. De même
$$\int _0^{\pi/4}\sin (2x)\,\mathrm{d}x=\left[\frac{-\cos (2x)}{2}\right]_0^{\pi/4}=\frac{1}{2}\quad \mbox{et}\quad \int _0^{\pi/4}\cos (2x)\,\mathrm{d}x=\left[\frac{\sin (2x)}{2}\right]_0^{\pi/4}=\frac{1}{2},$$
$$\int _0^{1}\frac{\mathrm{d}x}{x+1}=\Big[\ln (x+1)\Big]_0^1=\ln 2\quad \mbox{et}\quad \int _0^{1}\frac{\mathrm{d}x}{x+2}=\Big[\ln (x+2)\Big]_0^1=\ln 3-\ln 2.$$
Enfin,
$$\int _0^{1}\frac{\mathrm{d}x}{(x+1)^2}=\left[\frac{-1}{x+1}\right]_0^1=\frac{1}{2}\quad \mbox{et}\quad \int _0^{1}\frac{\mathrm{d}x}{1+x^2}=\Big[\arctan x\Big]_0^1=\pi/4.$$
\end{explanations}
\end{question}
}
